%%%%%%%%%%  scan idx 110  |  folio 99  %%%%%%%%%%

\medskip
\noindent\textbf{Lema 4.4.}
Para todo esquema $X_0$ sobre $k$, existe una sucesión exacta
canónica
\[
0\longrightarrow H^1(X_0,\underline D^0_{X_0})
\longrightarrow[\underline M_{X_0}(k[\varepsilon])]
\longrightarrow H^0(X_0,\underline D^1_{X_0}),
\]
donde los $\underline D^i_{X_0}$ son los definidos en la sección 3.

\medskip
\noindent\textit{Demostración.}
Supongamos primero que $X_0$ es afín, y pongamos
$\Gamma(X_0,\mathcal O_{X_0})=A_0$. Como se observó en 4.1(b), la
categoría $\underline M_{X_0}(k[\varepsilon])$ es canónicamente
equivalente a una categoría $\underline E$ cuyos objetos son
$k[\varepsilon]$-álgebras planas $A$ dotadas de una aplicación de
$k[\varepsilon]$-álgebras $j_A:A\longrightarrow A_0$ que induce un
isomorfismo
\[
k\otimes_{k[\varepsilon]}A\sisi{\overset{\sim}{\relbar\joinrel\relbar\joinrel\relbar}}{\xrightarrow{\ \sim\ }}A_0.
\]
Equivalentemente, un objeto de $\underline E$ es una $k$-álgebra $A$
dotada de una sucesión exacta
\[
0\longrightarrow A_0\sisi{\overset{\varepsilon}{\relbar\joinrel\relbar\joinrel\relbar}}{\xrightarrow{\ \varepsilon\ }}A
\sisi{\overset{j_A}{\relbar\joinrel\relbar\joinrel\relbar}}{\xrightarrow{\ j_A\ }}A_0\longrightarrow0
\]
en la cual $\varepsilon(A_0)^2=0$ como ideal de $A$.

Una flecha $A_1\longrightarrow A_2$ en $\underline E$, en la primera
descripción, es una aplicación de $k[\varepsilon]$-álgebras
$\Phi:A_1\longrightarrow A_2$ tal que
$j_{A_1}=j_{A_2}\circ\Phi$. En la segunda descripción es una
aplicación de $k$-álgebras $\Phi:A_1\longrightarrow A_2$ que produce
un diagrama conmutativo
\[
\begin{tikzcd}[column sep=large,row sep=large]
0 \arrow[r]
& A_0 \arrow[r,"\varepsilon"] \arrow[d,equal]
& A_1 \arrow[r,"j_{A_1}"] \arrow[d,"\Phi"]
& A_0 \arrow[r] \arrow[d,equal]
& 0 \\
0 \arrow[r]
& A_0 \arrow[r]
& A_2 \arrow[r,"j_{A_2}"]
& A_0 \arrow[r]
& 0.
\end{tikzcd}
\]
Se sigue que, mediante $\Gamma(X_0,-)$,
$[\underline M_{X_0}(k[\varepsilon])]$ es canónicamente isomorfo al
conjunto de clases de isomorfía de extensiones de $k$-álgebras de
$A_0$ por $A_0$. En otras palabras, existe una biyección canónica
\[
[\underline M_{X_0}(k[\varepsilon])]
\sisi{\overset{\sim}{\relbar\joinrel\relbar\joinrel\relbar}}{\xrightarrow{\ \sim\ }}H^0(X_0,\underline D^1_{X_0})
\]
que evidentemente respeta las estructuras de espacio vectorial en
ambos lados.

%%%%%%%%%%  scan idx 111  |  folio 100  %%%%%%%%%%

En general, para todo abierto afín $U_0\subset X_0$, la restricción y
el caso afín dan la aplicación compuesta
\[
[\underline M_{X_0}(k[\varepsilon])]
\longrightarrow[\underline M_{U_0}(k[\varepsilon])]
\xrightarrow{\ \sim\ }H^0(U_0,\underline D^1_{X_0}).
\]
Estas aplicaciones dan una aplicación canónica
\[
[\underline M_{X_0}(k[\varepsilon])]
\longrightarrow H^0(X_0,\underline D^1_{X_0}).
\]
Su núcleo está formado por las clases de isomorfía de las
deformaciones localmente triviales de $X_0$ sobre
$k[\varepsilon]$, es decir, de las deformaciones localmente isomorfas
a $X_0\otimes_k k[\varepsilon]$. El haz de automorfismos de
$X_0\otimes_k k[\varepsilon]$ que inducen la identidad sobre $X_0$
es
\[
\underline D^0_{X_0}
=\underline{\operatorname{Der}}_k
  (\mathcal O_{X_0},\mathcal O_{X_0}).
\]
Las clases de isomorfía de tales deformaciones localmente triviales
se identifican, por tanto, canónicamente con
$H^1(X_0,\underline D^0_{X_0})$, lo que da la sucesión exacta
afirmada. \hfill$\square$
